

\newpage

\paragraph{Problem 1 answer}\GradescopeComment{This should be on page 3 (top) of your PDF.}

%%%%%%%%%%%%%%%%%%%%%%%%%%%%%%%%%%%%%%%%%%%%%%%%%%%%%% 
%%%%%%%%%%%%%%%%%%%%%%%%%%%%%%%%%%%%%%%%%%%%%%%%%%%%%% 
%% 
%% PROBLEM 1(a) ANSWER 
%% 
%%%%%%%%%%%%%%%%%%%%%%%%%%%%%%%%%%%%%%%%%%%%%%%%%%%%%% 
%%%%%%%%%%%%%%%%%%%%%%%%%%%%%%%%%%%%%%%%%%%%%%%%%%%%%% 

\paragraph{(a)}

% Define your reduction $r$
% by defining, for every possible instance $A[1..n, 1..m]$ of  \prob{Tiling by Bars},
% what output $r(A)$ the reduction gives.

Given an instance $A[1..n, 1..m]$ of \prob{Tiling by Bars},
the reduction $r$ outputs the pair $(G, k) = r(A)$ where
$G=(U,W,E)$ and $k$ are defined as follows:

\begin{blue}
  \REPLACEME
\end{blue}

%%%%%%%%%%%%%%%%%%%%%%%%%%%%%%%%%%%%%%%%%%%%%%%%%%%%%% 
%%%%%%%%%%%%%%%%%%%%%%%%%%%%%%%%%%%%%%%%%%%%%%%%%%%%%% 
%% 
%% PROBLEM 1(b) ANSWER 
%% 
%%%%%%%%%%%%%%%%%%%%%%%%%%%%%%%%%%%%%%%%%%%%%%%%%%%%%% 
%%%%%%%%%%%%%%%%%%%%%%%%%%%%%%%%%%%%%%%%%%%%%%%%%%%%%% 

\paragraph{(b)}  

% Show that your reduction is a \emph{polynomial-time} reduction,
% by explaining carefully how, given $A$, one can compute the output $r(A)$ in time polynomial in the input size, $n\, m$.

Given an instance $A[1..n, 1..m]$ of \prob{Tiling by Bars}, the output $(G, k)$
of the reduction defined in Part (a) can be computed in time polynomial in the size of $A$, that is, $n\,m$, as follows:

\begin{blue}
  \REPLACEME
\end{blue}

%%%%%%%%%%%%%%%%%%%%%%%%%%%%%%%%%%%%%%%%%%%%%%%%%%%%%% 
%%%%%%%%%%%%%%%%%%%%%%%%%%%%%%%%%%%%%%%%%%%%%%%%%%%%%% 
%% 
%% PROBLEM 1(c) ANSWER 
%% 
%%%%%%%%%%%%%%%%%%%%%%%%%%%%%%%%%%%%%%%%%%%%%%%%%%%%%% 
%%%%%%%%%%%%%%%%%%%%%%%%%%%%%%%%%%%%%%%%%%%%%%%%%%%%%% 

\paragraph{(c)}

% Show the ``if'' direction.
% Given an arbitrary \prob{Tiling by Bars} instance $A$,
% let $G=(U, W, E)$ and $k$ be the output of your reduction $r$ given $A$.  That is, $(G, k) = r(A)$.
% Describe carefully how, given any bipartite matching $A$ in $G$ of size at least $k$,
% one could obtain from $A$ a valid covering of $A$ by bars.

 Fix an arbitrary \textsc{Tiling by Bars} instance $A[1..n, 1..m]$. 
 % [review 2026-09-27] fix: K: reduction output written r(I) but the instance is A -> r(A)
 Let $G=(U,W,E)$ and $k$ be the output $r(A)$ of the reduction $r$ given $A$.
 Consider any matching $M$ in $G$ of size at least $k$. 
 From $M$, we can obtain a valid tiling $T$ (of $A$ by bars) as follows.
 
\begin{blue}
  \REPLACEME
\end{blue}

\newpage                        % please don't delete this, doing so will mess up the page layout

\paragraph{Problem 1 answer, continued}\GradescopeComment{This should be on page 4 (top) of your PDF.}

%%%%%%%%%%%%%%%%%%%%%%%%%%%%%%%%%%%%%%%%%%%%%%%%%%%%%% 
%%%%%%%%%%%%%%%%%%%%%%%%%%%%%%%%%%%%%%%%%%%%%%%%%%%%%% 
%% 
%% PROBLEM 1(d) ANSWER 
%% 
%%%%%%%%%%%%%%%%%%%%%%%%%%%%%%%%%%%%%%%%%%%%%%%%%%%%%% 
%%%%%%%%%%%%%%%%%%%%%%%%%%%%%%%%%%%%%%%%%%%%%%%%%%%%%% 

\paragraph{(d)}

% Now show the ``only if'' direction.
% Given an arbitrary \prob{Tiling by Bars} instance $A$,
% let $G=(U, W, E)$ and $k$ be the output of your reduction $r$ given $A$.
% Describe carefully how, given any valid covering of $A$ by bars,
% one could obtain a bipartite matching $M$ in $G$ of size at least $k$.

Fix an arbitrary \textsc{Tiling by Bars} instance $A[1..n, 1..m]$. 
% [review 2026-09-27] fix: K: reduction output written r(I) but the instance is A -> r(A)
Let $G=(U,W,E)$ and $k$ be the output $r(A)$ of the reduction $r$ given $A$.
Consider any valid tiling $T$ of $A$ by bars.
From the tiling $T$, we can obtain a matching $M$ of size at least $k$ as follows.

\begin{blue}
  \REPLACEME
\end{blue}

%%%%%%%%%%%%%%%%%%%%%%%%%%%%%%%%%%%%%%%%%%%%%%%%%%%%%% 
%%%%%%%%%%%%%%%%%%%%%%%%%%%%%%%%%%%%%%%%%%%%%%%%%%%%%% 
%% 
%% PROBLEM 1(e) ANSWER 
%% 
%%%%%%%%%%%%%%%%%%%%%%%%%%%%%%%%%%%%%%%%%%%%%%%%%%%%%% 
%%%%%%%%%%%%%%%%%%%%%%%%%%%%%%%%%%%%%%%%%%%%%%%%%%%%%% 

\paragraph{(e)}

% Give pseudo-code for an algorithm that uses $R$ and $B$ to solve the \prob{Tiling by Bars} decision problem
% in polynomial time.  Your algorithm may call $B$ and $R$ as subroutines.  You don't need to define $B$ or define $R$.
% Your pseudo-code should have just a few short lines. \emph{Hint: Section 8.2.4, Lemma 8.6.}

Here is the algorithm that uses $R$ and $B$ to solve \textsc{Tiling by Bars}:

\bigskip

\begin{steps}
\item[] \textsf{TileWithBars}$(A[1..n, 1..m])$:
  \begin{blue}
    \step \REPLACEME
    \step $\vdots$              % you might need more than two steps
  \end{blue}  
\end{steps}

%%%%%%%%%%%%%%%%%%%%%%%%%%%%%%%%%%%%%%%%%%%%%%%%%%%%%% 
%%%%%%%%%%%%%%%%%%%%%%%%%%%%%%%%%%%%%%%%%%%%%%%%%%%%%% 
%% 
%% PROBLEM 1 COLLABORATORS AND CITATIONS 
%% 
%%%%%%%%%%%%%%%%%%%%%%%%%%%%%%%%%%%%%%%%%%%%%%%%%%%%%% 
%%%%%%%%%%%%%%%%%%%%%%%%%%%%%%%%%%%%%%%%%%%%%%%%%%%%%% 

\vfill
\begin{answerBox}{2in}          % please don't edit this line, doing so will throw off the page layout

  \paragraph{Problem \arabic{problem} collaboration and citations}~\GradescopeComment
  {This should be on page 4 (bottom) of your PDF.}

  I collaborated on this problem with the following people:
  \begin{blue}
    \begin{itemize}
    \item \REPLACEME                    % write "none" if none 
    \end{itemize}
  \end{blue}

  My answers use ideas from the following sources (other than the lecture notes and textbooks): 
  \begin{blue}
    \begin{itemize}
    \item \REPLACEME                    % write "none" if none 
    \end{itemize}
  \end{blue}

\end{answerBox}


%%% Local Variables:
%%% mode: latex
%%% TeX-master: "homework_10"
%%% End:
